\documentclass[11pt]{article}

\usepackage[T1]{fontenc}
\usepackage[utf8]{inputenc}
\usepackage{lmodern}
\usepackage{microtype}
\usepackage{geometry}
\geometry{margin=1in}
\usepackage{setspace}
\onehalfspacing
\usepackage{amsmath,amssymb,amsthm}
\usepackage{mathtools}
\usepackage{physics}
\usepackage{bm}
\usepackage{graphicx}
\usepackage{tikz}
\usetikzlibrary{arrows.meta,calc,decorations.pathreplacing}
\usepackage{hyperref}
\usepackage[nameinlink,noabbrev]{cleveref}
\usepackage[numbers,sort&compress]{natbib}

\title{The Geometry of Spacetime:\\
A Temporal-Radial Hyperspherical Model of Cosmological Expansion}
\author{[Name]}
\date{Draft manuscript}

\begin{document}
\maketitle

\begin{abstract}
Modern cosmology denies the existence of a spatial center of the universe: in a homogeneous and isotropic universe, every observer sees themselves at the center of their own observable region. This paper develops a geometrical interpretation in which the apparent absence of a spatial center is reconciled with a meaningful notion of a temporal center. The model represents three-dimensional space at cosmic time $t$ as a hypersurface $\Sigma_t$ embedded as the boundary of a growing four-dimensional sphere. The radial direction of the embedding is not interpreted as an additional traversable spatial coordinate, but as a representation of cosmic time. The center of the construction is therefore not a place inside the universe, but the limiting zero-radius boundary corresponding to the beginning of cosmic time. The model is formulated in relation to the closed Friedmann--Lema\^itre--Robertson--Walker metric and used to derive past-directed null curves, showing why distant observation corresponds geometrically to looking inward toward earlier hypersurfaces rather than along the tangent of the present universe. The observable horizon, cosmic microwave background, accelerated expansion, and a possible angular-drift extension are discussed. The paper emphasizes that the temporal-radial hypersphere is mathematically useful as an embedding interpretation, but any claim beyond standard FLRW cosmology requires additional dynamical assumptions and observational tests.
\end{abstract}

\noindent\textbf{Keywords:} spacetime geometry; hyperspherical universe; FLRW cosmology; cosmic time; observable universe; Big Bang; null geodesics; accelerated expansion; cosmological horizon.

\section{Introduction}

The question ``Where is the center of the universe?'' appears simple, but in modern cosmology it is incorrectly framed. The observable universe is centered on the observer, not because the observer occupies a privileged physical location, but because observation is limited by the finite speed of light and the finite age of the universe. A distant observer billions of light-years away would also see themselves at the center of their own observable universe. The cosmological principle, according to which the universe is homogeneous and isotropic on sufficiently large scales, removes the possibility of a unique spatial center.

However, the absence of a spatial center does not necessarily remove the possibility of a geometrical origin. The Big Bang is not a point in pre-existing space from which matter exploded outward. Rather, it is the early boundary of spacetime itself. The proper scientific reformulation is therefore not: ``Where is the center of the universe?'' but: ``Can the beginning of cosmic time be represented geometrically as the temporal center of spacetime?''

This paper develops a model in which three-dimensional space is represented as the hypersurface of a growing four-dimensional sphere. The radial coordinate of the embedding is interpreted as cosmic time. In this view, the center of the sphere is not a location in the three-dimensional universe. It is the temporal origin: the limiting condition $A(t)\rightarrow0$, corresponding to the beginning of cosmic time.

The proposal is not that galaxies move through a higher-dimensional spatial volume toward or away from a central point. Instead, the universe at each cosmic time is represented as a spatial hypersurface, and the succession of these hypersurfaces forms the history of cosmic expansion. This is closely related to the closed FLRW model, but the emphasis here is interpretive: the ``inside'' of the sphere is not ordinary space. It is a geometrical representation of the past.

The standard cosmological model, $\Lambda$CDM, remains highly successful. Planck and later analyses remain broadly consistent with spatial flatness within uncertainties \citep{planck2018vi,tristram2024}.

\section{From a Spatial Center to a Temporal Center}

Let the universe at cosmic time $t$ be represented by a three-dimensional spacelike hypersurface, $\Sigma_t$. In this model, each $\Sigma_t$ is embedded in an auxiliary four-dimensional Euclidean space as the surface of a four-dimensional sphere:
\[
\Sigma_t \cong S^3(A(t)).
\]
The cosmic history is represented as a foliation,
\[
\mathcal{M} = \bigcup_t \Sigma_t,
\]
where each hypersurface corresponds to a different cosmic time.

The Big Bang corresponds to the limiting condition $A(t)\rightarrow0$. Thus the ``center'' of the construction is not a point on the present spatial hypersurface, nor a coordinate that an observer can travel to through ordinary space.

\[
\boxed{\text{The center of the universe is not a where. It is a when.}}
\]

\section{Relation to FLRW Cosmology}

For a positively curved FLRW geometry,
\[
\dd s^2=-c^2\dd t^2+A^2(t)\left[\dd\chi^2+\sin^2\chi\left(\dd\theta^2+\sin^2\theta\,\dd\phi^2\right)\right].
\]
Writing $A(t)=a(t)R_c$ separates the scale factor from curvature radius. With $a(t_0)=1$, the model is mathematically closest to a closed FLRW universe.

Observationally, current data favor near-flat geometry, so the model is best viewed either as (i) an embedding visualization, or (ii) a closed-universe scenario with $A(t_0)$ far larger than the observable horizon scale \citep{planck2018vi,tristram2024}.

\section{Temporal-Radial Embedding and the View Curve}

A point on the embedded three-sphere is
\[
\mathbf{X}(t,\chi,\theta,\phi)=A(t)\,\mathbf{n}(\chi,\theta,\phi),
\]
with $\mathbf{n}$ a unit vector on $S^3$. The embedding center $\mathbf{X}=0$ is not part of the present spatial slice; it appears only as the limit $A(t)\rightarrow0$.

For radial null propagation, $\dd s^2=0$ implies
\[
\frac{\dd\chi}{\dd t}=\pm\frac{c}{A(t)}.
\]
For light arriving at $t_0$,
\[
D_C(t)=\int_t^{t_0}\frac{c\,\dd t'}{a(t')}=c\int_0^z\frac{\dd z'}{H(z')}.
\]
This is the relevant observable distance, not the Euclidean arc-length in the auxiliary embedding space \citep{hogg1999}.

\section{CMB, Horizons, and Acceleration}

The cosmic microwave background can be represented as the intersection
\[
\mathrm{CMB}=J^-(p_0)\cap\Sigma_{t_{\mathrm{LSS}}},
\]
of the observer's past light cone with the last-scattering hypersurface. The particle horizon at the present epoch is
\[
D_{\mathrm{hor}}(t_0)=\int_0^{t_0}\frac{c\,\dd t}{a(t)}.
\]
In standard cosmology this is approximately $46$--$47$ billion light-years in comoving distance.

Expansion and acceleration satisfy
\[
H(t)=\frac{\dot A}{A},\qquad \ddot A>0\,\,\text{(acceleration)}.
\]
The underlying dynamics remain those of the Friedmann equations; the embedding view is interpretive, not a replacement for dark energy physics \citep{riess1998,perlmutter1999}.

\section{Speculative Extension: Angular Drift}

A phenomenological drift extension can be written
\[
\dd s^2=-c^2\dd t^2+A^2(t)\left[\dd\chi^2+\sin^2\chi\left(\dd\theta^2+\sin^2\theta(\dd\phi-\Omega(t,\chi)\dd t)^2\right)\right].
\]
A photon would accumulate
\[
\Delta\phi\approx\int_{t_e}^{t_0}\Omega(t,\chi(t))\,\dd t.
\]
In the baseline model $\Omega=0$. Nonzero drift is speculative and strongly constrained by isotropy/polarization tests \citep{planck2018vii,gruppuso2020}.

\section{Black Holes and Local Temporal Deformation}

Within general relativity, black holes represent local causal structure deformation, not paths to the cosmological temporal origin. The ``icicle'' metaphor can be used qualitatively for severe time-dilation effects, but should not be interpreted literally as travel to the Big Bang.

\section{Observational Constraints and Falsifiability}

As an interpretation of FLRW, the model introduces no new predictions. To become a distinct theory, it must introduce measurable effects, such as detectable positive curvature, matched-circle signatures, nonzero angular drift, or a modified expansion history distinguishable from $\Lambda$CDM.

\section{Conclusion}

The model reframes the center problem: spacetime has no spatial center in the FLRW sense, but it does possess a temporal origin. In this interpretation, the Big Bang is the zero-radius boundary of cosmic time, and distant observation corresponds to traversing null curves toward earlier hypersurfaces.


\bibliographystyle{unsrtnat}
\bibliography{references}

\end{document}
